% Member 4 - final report, III.A Problem Formulation and Notation
\subsection{Problem Formulation and Notation}
\label{sec:notation}

We study single-source domain generalization. The labelled data comes from
one source domain, $D_s$, and the model is evaluated on target domains
$D_t$ that are never seen during training. In our setting $D_s$ is the
Photo domain of PACS, and the three targets are Art Painting, Cartoon and
Sketch. A training image $x \in D_s$ carries a label $y$ from a set of
seven object classes, and a classifier $f_\theta$ must predict $y$ for an
image from any of the four domains.

Style augmentation adds to $D_s$ a set of generated images. For each
training image $x$ and a style reference $s$ drawn from a frozen pool of
external paintings, a generator produces a stylised image $g(x, s)$. The
augmented training set is
\begin{equation}
D_s^{\mathrm{aug}} = D_s \;\cup\; \{ g(x, \sigma(x)) : x \in D_s \},
\label{eq:augset}
\end{equation}
where $\sigma$ is a fixed assignment of a style reference to each content
image, chosen once and frozen. The pairing $\sigma$ is not arbitrary: it is
drawn so that every style in the pool is used, and so that the assignment
does not correlate the style with the class. Both properties are checked
mechanically by the fairness checker before any model is trained
(Section~\ref{sec:fairness}).

The stylised image is not an arbitrary function of $(x, s)$. It is produced
by a feature-level blend of the two images, and the blend is governed by a
\emph{strength field} $\alpha(x, s) \in [0,1]^{H \times W}$, defined on the
spatial grid of the feature map. Writing $F_c$ for the content feature map
and $F_s$ for the style feature map, the generator computes
\begin{equation}
F = \alpha(x,s) \odot \mathrm{AdaIN}(F_c, F_s) + \bigl(1 - \alpha(x,s)\bigr) \odot F_c ,
\label{eq:blend}
\end{equation}
and decodes $F$ into the output image. Equation~\eqref{eq:blend} is the
single point at which our method and the AdaIN baseline differ: the
baseline uses a constant $\alpha$, ours uses a field.

This formulation makes the two quantities the study separates explicit.
The \emph{mean strength} of a configuration is
\begin{equation}
\bar{\alpha} = \mathbb{E}_{x}\left[\frac{1}{HW}\sum_{h,w} \alpha(x,s)(h,w)\right],
\label{eq:meanstrength}
\end{equation}
and the \emph{spatial distribution} is everything about $\alpha(x,s)$ beyond
its mean. Equation~\eqref{eq:meanstrength} is the quantity that turns out
to predict the accuracy: in our experiments the cross-domain accuracy is a
function of $\bar{\alpha}$ across configurations, and the spatial
distribution of the strength is not measurably influential once $\bar{\alpha}$
is held fixed (Section~\ref{sec:strength_energy}). Reporting $\bar{\alpha}$
for every configuration is therefore a precondition for the comparisons in
this paper to be interpretable, and we report it for each one.

Two degenerate cases of $\alpha$ anchor the design. $\alpha \equiv 0$ gives
the encoder--decoder reconstruction without any style. $\alpha \equiv c$ for
a constant $c$ gives plain AdaIN at strength $c$. Our method occupies the
space between them, and the experiments are designed so that each point of
that space can be compared with a uniform alternative of exactly the same
mean strength.
